\documentclass[10pt]{article}
\usepackage[margin=.85in]{geometry}
\usepackage{amsmath,amssymb,amsthm,booktabs,longtable}
\newcommand{\E}{\mathbb E}\newcommand{\Pp}{\mathbb P}
\newcommand{\CE}{C_{\rm E}}\newcommand{\R}{\mathbb R}
\newtheorem{theorem}{Theorem}\newtheorem{lemma}{Lemma}
\title{Elementary contraction of the thirty fixed-target entry panels}
\author{}\date{28 October 2026}
\begin{document}\maketitle
\noindent This is a conditional consequence of the accepted object
\texttt{INITIAL42\_FIXED\_TARGET\_COARSE\_V7}. It does not prove a count
comparison, confinement to half amplitude, or the finite base. The numerical
work is exact rational arithmetic on thirty probability intervals, with no
Fourier-kernel execution.

\paragraph{Inputs and notation.}
Let \(X\) be a centered variance-one fixed-target maximizing witness at
\(n\ge N=42\), with nonnegative excess over \(c=\CE\).
The previously established entry supplies \(nD\le ND_I\), the moment
ceilings in each closed probability panel \(I=[l,u]\), and an actual central
gap \((-A,B)\), where \(A,B>0\).
Here \(\rho=\E|X|^3\), \(\gamma=\E X^3\), \(\beta=\E[X|X|]\),
\(D=(3+\sqrt{10})\rho-\gamma-3J\), and \(J\) is the weighted gap sum.
The final active rho rows also give \(|X|\le R=157/50\):
this is the maximum of their individually acquired support radii.
All thirty nonempty panels, indexed 24 through 53, are retained initially.
Their data are bound by SHA256
\texttt{61b7ad43e20046d7cbf7157e3feb61696331731e7be418f6a5fa535c2eee8bae}.

Write \(p=\Pp(X>0)\), \(q=1-p\), \(m=\E X_+=\E X_->0\).
The conditional means are \(b=m/p\), \(a=m/q\), and
\[
 \theta=1-\frac{m^2}{pq},\qquad
 \ell=\frac m{pq}=\sqrt{\frac{1-\theta}{pq}},\qquad
 \delta=1-\frac{A+B}{\ell},\qquad
 \lambda=\frac{n\theta}{\ell^2}.
\]
Thus \(\delta\) is a count-coordinate inward displacement, not \(D\).
Let \(V=b\) on the positive sign and \(V=-a\) on the negative sign;
write \(W=X-V\). Define the two measures on nonnegative magnitudes by
\(\int f(x)d\mu_\pm(x)=\E[W^2\mathbf1_{\{\pm X>0\}}f(|X|)]\).
Their masses sum to \(\theta\); the central gap ensures there is no atom at zero.
Also \(A\le a\), \(B\le b\).

\paragraph{Retained analytic premise.}
We use the actual-law outer-gap inequality
\[
 D\ge(\sigma-1)\E X_+^3+(\sigma+1)\E X_-^3
       -6mAB+6BR_++6AR_-,
 \quad \sigma=\sqrt{10},\quad
 R_+=\E X_+^2-mB,\ R_-=\E X_-^2-mA.
\]
It is the input proved in the earlier weighted-gap analysis; arithmetic
evaluation does not prove it. Substituting
\(\E X_+^3=mb^2+\int(x+2b)d\mu_+\) and
\(R_+=\mu_+(\R)+m(b-B)\), and their negative-sign counterparts, gives
\begin{equation}\label{budget}
 D\ge B_0+E_0+\int k_+d\mu_++\int k_-d\mu_-,
\end{equation}
where, with \(p_*=(4-\sqrt{10})/2\),
\[
\begin{split}
 B_0&=\frac{(6+2\sigma)(1-\theta)^{3/2}}{\sqrt{pq}}(p-p_*)^2,\\
 E_0&=6m\{(A+B)(a-A+b-B)+(a-A)(b-B)\},\\
 k_+(x)&=(\sigma-1)(x+2b)+6B,\qquad
 k_-(x)=(\sigma+1)(x+2a)+6A .
\end{split}
\]
The deterministic identity is
\((\sigma+1)p^2+(\sigma-1)q^2-6pq=(6+2\sigma)(p-p_*)^2\).
Every displayed cost is nonnegative.

\paragraph{A whole-panel iteration.}
Set \(v_-=\min_{p\in I}pq\), \(v_+=\max_{p\in I}pq\).
The maximum includes \(1/4\) if \(I\) crosses \(1/2\).
Use \(\sigma_l=3.162277660<\sigma<3.162277661=\sigma_u\).
Initially \(D\le D_I<\sigma_l-1\), and the gap bound
\(J\le\rho+2mAB\) gives
\[
 AB\ge z:=(\sigma_l-1-D_I)/3>0.
\]
Indeed \(\rho\ge1,\gamma\le\rho,m\le1/2\) imply this inequality.
With \(a_u=\sqrt{u/(1-u)}\), \(b_u=\sqrt{(1-l)/l}\), obtain
\[
 A\ge A_l=z/b_u,\quad B\ge B_l=z/a_u,\quad
 (A+B)/\ell\ge g=2\sqrt{zv_-},\quad \delta\le1-g.
\]
For a currently proved \(\theta\le\Theta\), lower-bound the binary charge by
\[
 b_I=\frac{(6+2\sigma_l)(1-\Theta)^{3/2}}{\sqrt{v_+}}
        \operatorname{dist}(I,[ (4-\sigma_u)/2,(4-\sigma_l)/2])^2.
\]
All square roots are directed outward. If \(r=D_I-b_I<0\), the entire
closed panel is impossible. Otherwise set
\[
 a_l=\sqrt{(1-\Theta)l/(1-l)},\quad
 b_l=\sqrt{(1-\Theta)(1-u)/u},
\]
\[
 K=\min\{(\sigma_l-1)(B_l+2b_l)+6B_l,\
          (\sigma_l+1)(A_l+2a_l)+6A_l\}>0.
\]
Equation~\eqref{budget} gives \(\theta\le r/K\).
Intersect this bound with \(\Theta\), then recompute \(b_I,r\).
Put \(m_l=\sqrt{(1-\Theta)v_-}\) and
\(\ell_l=\sqrt{(1-\Theta)/v_+}\). Since
\(E_0\ge6m\ell^2g\delta\), we may update
\[
 \delta\le\min\{\delta_{\rm old},r/(6m_l\ell_l^2g)\},\quad
 g\ge1-\delta,\quad
 A_l\gets\max\{A_l,\ell_l(l-\delta)_+\},\quad
 B_l\gets\max\{B_l,\ell_l(1-u-\delta)_+\}.
\]
Only previously proved bounds enter each update. Zero residual forces
zero noise and displacement; no division by that residual is needed.
The implementation performs eight rounds, rounding upper bounds upward and
lower bounds downward to \(10^{-10}\).

\paragraph{Moment and weighted-noise exports.}
For every orientation, including \(p>1/2\),
\begin{equation}\label{beta}
 1-2u-2(1-u)\Theta\le\beta\le1-2l+2l\Theta.
\end{equation}
This follows from
\(\beta=(1-\theta)(1-2p)+\mu_+(\R)-\mu_-(\R)\).
The formerly convenient upper replacement \(1-2l+\Theta\) requires
\(l\le1/2\) and must not be continued past that point.

For a coherent third-moment cap, put
\[
 L=R+2\max(a_u,b_u),\quad
 k=\min\{\sigma_l-1+6B_l/(R+2b_u),\
          \sigma_l+1+6A_l/(R+2a_u)\}.
\]
The exact identity
\(\rho=\rho_B(p)(1-\theta)^{3/2}+\sum_\pm\int(x+2a_\pm)d\mu_\pm\)
gives a noise charge at most \(\min(L\theta,r/k)\).
Let \(r_B=(1-2v_-)/\sqrt{v_-}\) and
\(s\le(3/2)\sqrt{1-\Theta}\). Integrating the derivative proves
\((1-\theta)^{3/2}\le1-s\theta\) for \(0\le\theta\le\Theta\).
The maximum of this affine, piecewise expression is therefore bounded by
\[
 \rho\le r_B+(L-sr_B)_+\min\{\Theta,r/(kL)\}.
\]
Intersect with the prior rho cap.

Finally \eqref{budget} implies the joint inequality
\[
 D\ge b_I+6m(A_l+B_l)\ell\delta+K\theta.
\]
Multiplying by the actual \(n\), using \(nD\le ND_I\),
\(m\ell=1-\theta\), and \(n\theta=(1-\theta)\lambda/(pq)\),
proves, when \(r>0\),
\[
 \frac{\delta}{r/[6(1-\Theta)(A_l+B_l)]}
 +\frac{\lambda}{Nv_+r/[K(1-\Theta)]}\le1.
\]
This retains the weighted resource; it is not obtained by multiplying an
arbitrary unweighted noise cap by \(N\).

\begin{theorem}[Conditional all-panel result]
The exact eight-round calculation excludes panels 24 and 53.
On all 28 remaining panels,
\[
 \frac14\le p\le\frac{53}{100},\quad
 \theta\le.0594425514,\quad \delta\le.0719539227,\quad
 \rho\le1.4547288545,\quad \lambda\le.6314459997,
\]
\[
 1.957063175\le\ell\le2.3094010768.
\]
All endpoints and both signs are included.
\end{theorem}
\begin{proof}
The previous inequalities justify each recorded rational update.
The input is the full accepted closed cover, and only a strictly negative
residual excludes a panel. All 28 retained outputs and every update are in
the exact log; the stated constants are their outward maxima or minima.
No sampled parameter or favorable moment optimizer is used.
\end{proof}

\paragraph{A direct support-strip check and its limit.}
The retained fixed-target contact theorem, with sharp predecessor and
nonnegative excess, bounds a support contact by \(H_n(H,\pm t)+cQ(t)\),
where
\(Q(t)=-t^3+3rt^2/2+3bt+(8/875)r\).
Here \(r\) is a rho upper bound, and \(b=\beta_u\) on the positive ray,
\(b=-\beta_l\) on the negative ray.
For \(t>\sqrt3\), set
\[
 A_t=(t^3-3t)/6,\quad H_4=t^4-6t^2+3,\quad S=81/80,
\quad E_t=t^5/4100+3t/820+17(3t^2-2)/78720.
\]
The heat expansion used in the preceding entry proof remains valid here:
its fifth-derivative error has coefficient \(3t^2-2\), since \(t^2-1>0\).
Using \(\|\phi'''\|_\infty<2\phi(0)\), \(\phi(0)<2/5\) and
\(\sqrt{n-1}>6\), without any sign assumption on \(H_4\), yields
\[
 H_n(H,\pm t)\le \tfrac25S^3(A_t+|H_4|/72)+S^3E_t.
\]
Thus a negative boundary value of this expression plus
\((819/2000)Q(t)\) excludes that boundary, provided \(Q(t)<0\).
The exact derivative estimate is
\(|\partial_tH_n|<\tfrac25(t^2/2+1/2+8/875)\).
Check \(Q,Q',Q''<0\) and
\[
 L=\tfrac25(t^2/2+1/2+8/875)+(819/2000)Q',\quad L',L''<0.
\]
Their higher derivatives are negative constants, so both rays are excluded
beyond the boundary, for every raw threshold \(H\).
These Gaussian/contact premises are retained analytical inputs; none is
proved by a numerical PASS.

Choose the two boundaries
\(t_-=\ell_l(l+3/5)\) and \(t_+=\ell_l(1-u+3/5)\).
All satisfy \(t_\pm>\sqrt3\). Both pass on panels 40--52, proving
\(Y=X/\ell+p\in(-3/5,8/5)\) there.
On panels 25--39 the negative ray passes but the positive boundary
upper estimate is nonnegative. This is failure of that sufficient estimate,
not a counterexample to the strip theorem. No strip claim is made on those
fifteen panels, and no half-amplitude confinement is claimed anywhere.

\paragraph{Adversarial examples.}
The standardized three-point law with atoms and masses
\[
 (-26/25,12/25),\quad(1/2,52/6977),\quad
 (7226/7475,89401/174425)
\]
has \(p=13/25,\theta=1/625,\beta=-599/15625\).
Its beta exceeds \(1-2p+\theta=-24/625\) by \(1/15625\).
It verifies the general formula~\eqref{beta} and refutes the half-interval
shortcut outside its hypothesis; it is not a counterexample to Berry--Esseen.

A standardized Bernoulli law at \(p=13/50\) has
\(\theta=\delta=0\), \(\beta=12/25\),
\[
 \rho=\frac{769}{25\sqrt{481}},\qquad
 D=\frac{769\sqrt{10}-2043}{25\sqrt{481}}.
\]
It satisfies the remaining moment/defect constraints of panel 25.
Hence these necessary constraints alone cannot exclude that whole panel.
This law is not asserted to be a nonnegative-excess maximizing witness.
The checker also verifies the gap-cost identity on symmetric and skew
Bernoulli laws, their reflection, a four-point law and small distant atoms.
Seven invalid controls reject. All work remains separate from the finite base.
\end{document}